\subsection{対称性跳躍の全環と全次数の検証法}

$B_{1,1}$ の対称性保持環は完全に決まり、
\[
 R_Z=\mathbf C[[u_1,\ldots,u_5,\varepsilon]]/(\varepsilon^2)
\]
である。従って前版の $A[[t]]/(t(t-b),tJ)$ という表示の未知量は
$b=J=0$ となる。被約部分は前版の五変数 $\mathbf G_a$ 族である。
全体は次元五の局所完全交叉、Gorenstein、重複度二で、埋込成分を持たない。

再利用すべき判定は次のものである。
埋込次元 $k+1$ の完備局所環 $R$ に非零な二次関係があり、
$\mathbf C[[u_1,\ldots,u_k,\varepsilon]]/(\varepsilon^2)$ 上の族へ
接空間同型を与える写像があれば、$R$ 自体がこの平方零環に同型となる。
最小表示からの合成を形式的逆関数定理で標準化すると、
定義イデアルは $(\varepsilon^2)$ に含まれる。
非零な二次関係は $\varepsilon^2$ の単元倍なので逆包含も成立する。
有限jetだけの一致からこの結論を推論しない。

今回の族は、既存の五変数族を $\mu_u,D_u$ とし、
\[
 \widetilde\mu=\mu_u+\varepsilon\Phi/\Delta,
 \mathcal D=D_u+\varepsilon\Psi/\Delta,
 \quad\varepsilon^2=0,\quad\Delta(0)=1
\]
とする有理式で構成した。
原点で $\Phi=U_6$、$\Psi=E_6$ となる。
二cochainの追加補正は原始射の右合成に支持を持つ二十二座標、
一次cochainの追加補正は百座標だけである。
分母は全次数十二、分子は十四以下であり、
元のCGG積から直接再構成した九百七十二個の結合則と
千五百八十四個の導分方程式が、分母除去後に多項式恒等式として全て零となる。
この直接検証は、構成時に選んだ百二十二行の線形方程式から独立である。
全有理係数は \texttt{vk\_double\_symmetry\_v39\_witness.json}、
検証コードは \texttt{verify\_vk\_double\_symmetry\_v39.py} に保持する。

$D_u^{12}=0$ と $\varepsilon^2=0$ から、補正 $\Psi$ の形によらず
$\mathcal D^{24}=0$ が従う。
一般に $D^n=0$ かつ $\varepsilon^2=0$ なら
$(D+\varepsilon E)^{2n}=0$ である。
一個の $\varepsilon$ を含む項は $D^aED^b$、$a+b=2n-1$ だからである。
したがって全非被約基底上で有限和の指数写像による $\mathbf G_a$ 作用を得る。

通常の三十五変数hullの形式座標と残余列の基底も選び直せる。
\[
 R_B=\mathbf C[[u_1,\ldots,u_5,\varepsilon,z_1,\ldots,z_{29}]],
 \quad I_B=(z_1,\ldots,z_{29},\varepsilon^2),
 \quad a=(z_1,\ldots,z_{29},\varepsilon^2,0^5)^t.
\]
従って任意の形式弧に沿う対称性消失次数は正確に
\[
 \nu=\min\{\operatorname{ord}z_1,\ldots,
                  \operatorname{ord}z_{29},2\operatorname{ord}\varepsilon\}
\]
であり、有限なら相対 $HH^2$ の捩れは $\mathcal O/(s^\nu)$ である。
無限大の弧は被約部分 $Y$ を通る。
また $\operatorname{Ext}_{R_B}^1(HH^2_{R_B},R_B)=R_B/I_B$ によって、
二重構造を相対Hochschild加群の双対から回復できる。

純粋三変数部分の対称性保持曲線を $\mathbf C[[\gamma]]$ とすると、
上の同型はこの環上で
$R_Z=\mathbf C[[u_1,\ldots,u_5,\gamma]]/(\gamma^2)$ と取れる。
理由は、被約五変数族が純粋部分を固定するので $\gamma=\varepsilon h(u)$ となり、
純粋部分への接写像が非零であるため $h$ が単元になることである。
$B_{0,3}$ の四変数対称性保持部分は純粋曲線の原点へ全次数で写る。
純粋部分では滑らかな曲線が存在するが、原始射の接合がその方向を
原点または平方零近傍へ切り詰める、という形で二つの共鳴の違いを説明できる。

標準六項補完の像も同時に直線化できる。
純粋部分への制限行列に補完接線を代入すると
$(229/120,117/20,-1571/60)^t$ となり、純粋 $\gamma$ 曲線の接線と一致する。
六項hullは $\mathbb C[[\gamma]]/(\gamma^2)$ なので、
この環への像は $u_i=a_i\gamma$ となる。
$u_i$ を $u_i-a_i\gamma$ へ変えれば、標準補完像のイデアルは
$(u_1,\ldots,u_5,z_1,\ldots,z_{29},\gamma^2)$ になる。
被約五次元部分は $\gamma=z_1=\cdots=z_{29}=0$ のままである。
従って標準補完像は、全対称性保持部分の横断的な長さ二の切片である。
この同時正規形は座標の選択に依存し、別の例外列を許す補完の分類ではない。

次に調べる価値があるのは、この平方零の厚みを一般の重み・共鳴へ拡張することと、
純粋三変数部分の対称性保持曲線を原始射の接合がどの次数で切り詰めるかという問題である。
今回の厚み二は全次数で確定したものであり、既存の外側の
$\mathbf C[[\gamma]]/(\gamma^{j+1})$ と同じ次数が一般に続くかは未決である。
